\documentclass[UTF8, 12pt, fontset=fandol, oneside]{ctexart}
\usepackage{researchnotes}

\title{第八章\quad 二次型}
\author{}
\date{}

\begin{document}

\maketitle

\section*{§8.1\quad 二次型的化简与矩阵的合同}

在解析几何中，我们曾经学过二次曲线及二次曲面的分类. 以平面二次曲线为例，一条二次曲线可以由一个二元二次方程给出：
\begin{equation}
ax^2+bxy+cy^2+dx+ey+f=0.
\tag{8.1.1}
\end{equation}
要区分 (8.1.1) 式是哪一种曲线 (椭圆、双曲线、抛物线或其退化形式)，我们通常分两步来做：首先将坐标轴旋转一个角度以消去 $xy$ 项，再作坐标轴的平移以消去一次项. 这里的关键是消去 $xy$ 项，通常的坐标变换公式为
\begin{equation}
\left\{\begin{array}{l}
x=x'\cos\theta-y'\sin\theta,\\
y=x'\sin\theta+y'\cos\theta.
\end{array}\right.
\tag{8.1.2}
\end{equation}
从线性空间与线性变换的角度来看，(8.1.2) 式表示平面上的一个线性变换，因此二次曲线分类的关键是给出一个线性变换，使 (8.1.1) 式中的二次项只含平方项. 这种情形也在空间二次曲面的分类时出现. 类似的问题在数学的其他分支、物理、力学中也会遇到. 为了讨论问题的方便，只考虑二次齐次多项式.

\noindent\textbf{定义 8.1.1}\quad 设 $f$ 是数域 $\mathbb K$ 上的 $n$ 元二次齐次多项式：
\begin{equation}
\begin{aligned}
f(x_1,x_2,\cdots,x_n)
={}&a_{11}x_1^2+2a_{12}x_1x_2+\cdots+2a_{1n}x_1x_n\\
&+a_{22}x_2^2+\cdots+2a_{2n}x_2x_n+\cdots+a_{nn}x_n^2,
\end{aligned}
\tag{8.1.3}
\end{equation}
称 $f$ 为数域 $\mathbb K$ 上的 $n$ 元二次型，简称二次型.

这里非平方项的系数采用 $2a_{ij}$ 主要是为了以后矩阵表示的方便.

\noindent\textbf{例 8.1.1}\quad 下列多项式都是二次型：
\[
f(x,y)=\frac12x^2+3xy+4y^2,
\]
\[
f(x,y,z)=2x^2-y^2+4z^2-2xy+(\sqrt3+2)xz,
\]
\[
f(x_1,x_2,x_3,x_4)=x_1^2+\sqrt5x_2^2-2x_3^2.
\]
而下列多项式都不是二次型：
\[
f(x,y)=x^2+3xy+y^2-2x+1,
\]
\[
f(x,y,z)=2x^2-y^2-2xy+2xz+3,
\]
\[
f(x_1,x_2,x_3,x_4)=x_1^3+x_2^2-2x_3^2+5x_1x_4.
\]

我们现在要用矩阵为工具来处理二次型. 用矩阵的乘法我们可以把 (8.1.3) 式写成矩阵相乘的形式：
\begin{equation}
f(x_1,x_2,\cdots,x_n)=x' Ax,
\tag{8.1.4}
\end{equation}
其中
\[
A=\left(\begin{array}{cccc}
a_{11}&a_{12}&\cdots&a_{1n}\\
a_{21}&a_{22}&\cdots&a_{2n}\\
\vdots&\vdots&&\vdots\\
a_{n1}&a_{n2}&\cdots&a_{nn}
\end{array}\right),\quad
x=\left(\begin{array}{c}
x_1\\
x_2\\
\vdots\\
x_n
\end{array}\right).
\]
在矩阵 $A$ 中，$a_{ij}=a_{ji}$ 对一切 $i,j$ 成立，也就是说矩阵 $A$ 是一个对称阵. 由此可知，给定数域 $\mathbb K$ 上的一个 $n$ 元二次型，我们就得到了 $\mathbb K$ 上的一个 $n$ 阶对称阵 $A$，称为该二次型的相伴矩阵或系数矩阵. 反过来，若给定 $\mathbb K$ 上的一个 $n$ 阶对称阵 $A$，则由 (8.1.4) 式，我们可以得到 $\mathbb K$ 上的一个二次型，称为对称阵 $A$ 的相伴二次型. 现在的问题是：用这样的方法得到的对称阵是否和二次型一一对应？回答这一点并不难. 设 $f=x'Ax=x'Bx$，我们要证明 $A=B$. 这等价于证明下面的结论：设 $A=(a_{ij})$ 是 $n$ 阶对称阵，若 $\alpha'A\alpha=0$ 对一切 $\alpha$ 成立，则 $A=O$. 令 $\alpha=e_i$ 是 $n$ 维标准单位列向量，则 $a_{ii}=e_i'Ae_i=0$. 再令 $\alpha=e_i+e_j(i\neq j)$，则
\[
0=(e_i+e_j)'A(e_i+e_j)=e_i'Ae_i+e_j'Ae_j+e_i'Ae_j+e_j'Ae_i=a_{ij}+a_{ji}.
\]
因为 $a_{ij}=a_{ji}$，故 $a_{ij}=0(i\neq j)$，于是 $A=O$. 这表明用对称阵来表示二次型时，系数矩阵是唯一的. 事实上，如果我们不限制矩阵是对称阵，则系数矩阵将不唯一，这样会给用矩阵方法研究二次型带来困难.

\noindent\textbf{例 8.1.2}\quad 求二次型 $f(x,y)=x^2+3xy+y^2$ 的矩阵表示.

\noindent\textbf{解}\quad
\[
f(x,y)=(x,y)\left(\begin{array}{cc}
1&\dfrac32\\
\dfrac32&1
\end{array}\right)
\left(\begin{array}{c}
x\\
y
\end{array}\right).
\]

\noindent\textbf{例 8.1.3}\quad 求二次型 $f(x_1,x_2,x_3,x_4)=x_1^2-x_2^2+2x_3^2+4x_4^2$ 相伴的对称阵.

\noindent\textbf{解}\quad 所求的对称阵为
\[
\left(\begin{array}{cccc}
1&0&0&0\\
0&-1&0&0\\
0&0&2&0\\
0&0&0&4
\end{array}\right),
\]
这是一个对角阵. 显然，一个 $n$ 元二次型只含平方项当且仅当它的相伴矩阵是一个 $n$ 阶对角阵.

\noindent\textbf{例 8.1.4}\quad 求对称阵
\[
\left(\begin{array}{ccc}
1&\dfrac12&-\sqrt2\\
\dfrac12&-1&0\\
-\sqrt2&0&2
\end{array}\right)
\]
所对应的二次型.

\noindent\textbf{解}\quad 所求之二次型为
\[
f(x_1,x_2,x_3)=x_1^2-x_2^2+2x_3^2+x_1x_2-2\sqrt2x_1x_3.
\]

二次型理论的基本问题是要寻找一个线性变换把它变成只含平方项的标准型. 由上面我们知道，二次型与对称阵一一对应，而线性变换可以用矩阵来表示. 自然地，二次型的变换与矩阵有着密切的关系，现在我们来探讨这个关系.

设 $V$ 是 $n$ 维线性空间，二次型 $f(x_1,x_2,\cdots,x_n)$ 可以看成是 $V$ 上的二次函数. 即若设 $V$ 的一组基为 $\{e_1,e_2,\cdots,e_n\}$，向量 $\alpha$ 在这组基下的坐标向量为 $x=(x_1,x_2,\cdots,x_n)'$，则 $f$ 便是向量 $\alpha$ 的函数. 现假设 $\{f_1,f_2,\cdots,f_n\}$ 是 $V$ 的另一组基，向量 $\alpha$ 在 $\{f_1,f_2,\cdots,f_n\}$ 下的坐标向量为 $y=(y_1,y_2,\cdots,y_n)'$. 记 $C=(c_{ij})$ 是从基 $\{e_1,e_2,\cdots,e_n\}$ 到基 $\{f_1,f_2,\cdots,f_n\}$ 的过渡矩阵，则
\[
\left(\begin{array}{c}
x_1\\
x_2\\
\vdots\\
x_n
\end{array}\right)
=
\left(\begin{array}{cccc}
c_{11}&c_{12}&\cdots&c_{1n}\\
c_{21}&c_{22}&\cdots&c_{2n}\\
\vdots&\vdots&&\vdots\\
c_{n1}&c_{n2}&\cdots&c_{nn}
\end{array}\right)
\left(\begin{array}{c}
y_1\\
y_2\\
\vdots\\
y_n
\end{array}\right),
\]
或简记为
\[
x=Cy.
\]
将上式代入 $f(x_1,x_2,\cdots,x_n)=x'Ax$，得
\[
f(x_1,x_2,\cdots,x_n)=y'C'ACy.
\]
显然，$C'AC$ 仍是一个对称阵，故 $y'C'ACy$ 是以 $y_1,y_2,\cdots,y_n$ 为变元的二次型，记为 $g(y_1,y_2,\cdots,y_n)$. 由此我们可以看出：若二次型 $f(x_1,x_2,\cdots,x_n)$ 所对应的对称阵为 $A$，则经过变量代换之后得到的二次型 $g(y_1,y_2,\cdots,y_n)$ 所对应的对称阵为 $C'AC$.

\noindent\textbf{定义 8.1.2}\quad 设 $A,B$ 是数域 $\mathbb K$ 上的 $n$ 阶矩阵，若存在 $n$ 阶非异阵 $C$，使
\[
B=C'AC,
\]
则称 $B$ 与 $A$ 是合同的，或称 $B$ 与 $A$ 具有合同关系.

不难证明，合同关系是一个等价关系，即

(1) 任一矩阵 $A$ 与自己合同，因为 $A=I'A I$;

(2) 若 $B$ 与 $A$ 合同，则 $A$ 与 $B$ 合同. 这是因为若 $B=C'AC$，则 $A=(C')^{-1}BC^{-1}=(C^{-1})'BC^{-1}$;

(3) 若 $B$ 与 $A$ 合同，$D$ 与 $B$ 合同，则 $D$ 与 $A$ 合同. 事实上，若 $B=C'AC$，$D=H'BH$，则 $D=H'C'ACH=(CH)'A(CH)$.

因为一个二次型经变量代换后得到的二次型的相伴对称阵与原二次型的相伴对称阵是合同的，又因为只含平方项的二次型其相伴矩阵是一个对角阵 (见例 8.1.3)，所以，化二次型为只含平方项的标准型等价于对称阵 $A$ 寻找非异阵 $C$，使 $C'AC$ 是一个对角阵. 这一情形与矩阵相似关系颇为类似，在相似关系下我们希望找到一个非异阵 $P$，使 $P^{-1}AP$ 成为简单形式的矩阵 (如 Jordan 标准型). 现在我们要找一个非异阵 $C$，使 $C'AC$ 为对角阵. 因此，二次型化简的问题相当于寻找合同关系下的标准型. 我们能找到这样的矩阵 $C$ 吗？

首先我们来考察初等变换和矩阵合同的关系.

\noindent\textbf{引理 8.1.1}\quad 对称阵 $A$ 的下列变换都是合同变换：

(1) 对换 $A$ 的第 $i$ 行与第 $j$ 行，再对换第 $i$ 列与第 $j$ 列；

(2) 将非零常数 $k$ 乘以 $A$ 的第 $i$ 行，再将 $k$ 乘以第 $i$ 列；

(3) 将 $A$ 的第 $i$ 行乘以 $k$ 加到第 $j$ 行上，再将第 $i$ 列乘以 $k$ 加到第 $j$ 列上.

\noindent\textbf{证明}\quad 上述变换相当于将一个初等矩阵左乘以 $A$ 后再将这个初等矩阵的转置右乘之，因此是合同变换. $\square$

\noindent\textbf{引理 8.1.2}\quad 设 $A$ 是数域 $\mathbb K$ 上的非零对称阵，则必存在非异阵 $C$，使 $C'AC$ 的第 $(1,1)$ 元素不等于零.

\noindent\textbf{证明}\quad 若 $a_{11}=0$，而 $a_{ii}\neq0$，则将 $A$ 的第一行与第 $i$ 行对换，再将第一列与第 $i$ 列对换，得到的矩阵的第 $(1,1)$ 元素不为零. 根据上述引理，这样得到的矩阵和原矩阵合同.

若所有的 $a_{ii}=0(i=1,2,\cdots,n)$，设 $a_{ij}\neq0(i\neq j)$，将 $A$ 的第 $j$ 行加到第 $i$ 行上，再将第 $j$ 列加到第 $i$ 列上. 因为 $A$ 是对称阵，$a_{ij}=a_{ji}\neq0$，于是第 $(i,i)$ 元素是 $2a_{ij}\neq0$，再用前面的办法使第 $(1,1)$ 元素不等于零. 根据上述引理，这样得到的矩阵和原矩阵仍合同，这就证明了结论. $\square$

\noindent\textbf{定理 8.1.1}\quad 设 $A$ 是数域 $\mathbb K$ 上的 $n$ 阶对称阵，则必存在 $\mathbb K$ 上的 $n$ 阶非异阵 $C$，使 $C'AC$ 为对角阵.

\noindent\textbf{证明}\quad 由上述引理，不妨设 $A=(a_{ij})$ 中 $a_{11}\neq0$. 若 $a_{i1}\neq0$，则可将第一行乘以 $-a_{11}^{-1}a_{i1}$ 加到第 $i$ 行上，再将第一列乘以 $-a_{11}^{-1}a_{i1}$ 加到第 $i$ 列上. 由于 $a_{i1}=a_{1i}$，故得到的矩阵的第 $(1,i)$ 元素及第 $(i,1)$ 元素均等于零. 由引理 8.1.1 可知，新得到的矩阵与 $A$ 是合同的. 依次这样做下去，可把 $A$ 的第一行与第一列除 $a_{11}$ 外的元素都消去，于是 $A$ 合同于下列矩阵：
\[
\left(\begin{array}{ccccc}
a_{11}&0&0&\cdots&0\\
0&b_{22}&b_{23}&\cdots&b_{2n}\\
0&b_{32}&b_{33}&\cdots&b_{3n}\\
\vdots&\vdots&\vdots&&\vdots\\
0&b_{n2}&b_{n3}&\cdots&b_{nn}
\end{array}\right).
\]
上式右下角是一个 $n-1$ 阶对称阵，记为 $A_1$. 因此由归纳假设，存在 $n-1$ 阶非异阵 $D$，使 $D'A_1D$ 为对角阵，于是
\[
\left(\begin{array}{cc}
1&O\\
O&D'
\end{array}\right)
\left(\begin{array}{cc}
a_{11}&O\\
O&A_1
\end{array}\right)
\left(\begin{array}{cc}
1&O\\
O&D
\end{array}\right)
=
\left(\begin{array}{cc}
a_{11}&O\\
O&D'A_1D
\end{array}\right)
\]
是一个对角阵. 显然
\[
\left(\begin{array}{cc}
1&O\\
O&D'
\end{array}\right)
=
\left(\begin{array}{cc}
1&O\\
O&D
\end{array}\right)' ,
\]
因此 $A$ 合同于对角阵. $\square$

\section*{习题 8.1}

1. 证明下列关于分块对称阵 $A$ 的变换是合同变换：

(1) 对换 $A$ 的第 $i$ 分块行与第 $j$ 分块行，再对换第 $i$ 分块列与第 $j$ 分块列；

(2) 将非异阵 $M$ 左乘 $A$ 的第 $i$ 分块行，再将 $M'$ 右乘第 $i$ 分块列；

(3) 将矩阵 $M$ 左乘 $A$ 的第 $i$ 分块行加到第 $j$ 分块行上，再将 $M'$ 右乘第 $i$ 分块列加到第 $j$ 分块列上.

2. 设 $\operatorname{diag}\{A_1,A_2,\cdots,A_m\}$ 是分块对角阵，其中 $A_i$ 都是对称阵，求证：$\operatorname{diag}\{A_1,A_2,\cdots,A_m\}$ 合同于 $\operatorname{diag}\{A_{i_1},A_{i_2},\cdots,A_{i_m}\}$，其中 $A_{i_1},A_{i_2},\cdots,A_{i_m}$ 是 $A_1,A_2,\cdots,A_m$ 的一个排列.

3. 设可逆阵 $A$ 和 $B$ 合同，求证：$A^{-1}$ 和 $B^{-1}$ 也合同.

4. 设矩阵 $A$ 和 $B$ 合同，求证：$A^*$ 和 $B^*$ 也合同.

5. 设矩阵 $A_1$ 和 $B_1$ 合同，$A_2$ 和 $B_2$ 合同，证明下列两个分块矩阵也合同：
\[
\left(\begin{array}{cc}
A_1&O\\
O&A_2
\end{array}\right),\quad
\left(\begin{array}{cc}
B_1&O\\
O&B_2
\end{array}\right).
\]

6. 证明：秩等于 $r$ 的对称阵等于 $r$ 个秩为 1 的对称阵之和.

7. 设 $n$ 阶复对称阵 $A$ 的秩为 $r$，求证：$A$ 必可分解为 $A=T'T$，其中 $T$ 是秩为 $r$ 的 $n$ 阶矩阵.

8. 设 $A$ 是 $n$ 阶反对称阵，证明 $A$ 合同于下列形状的矩阵：
\[
\operatorname{diag}\{S,\cdots,S,0,\cdots,0\};
\]
其中
\[
S=\left(\begin{array}{cc}
0&1\\
-1&0
\end{array}\right).
\]
特别，反对称阵的秩总是偶数.

9. 求证：实反对称阵的行列式总是非负实数.

10. 求证：元素全是整数的反对称阵的行列式是某个整数的平方.

\section*{§8.2\quad 二次型的化简}

如何将二次型化为只含平方项的标准型？在这一节里，我们将介绍配方法和初等变换法这两种方法.

\subsection*{一、配方法}

配方法的基础是运用下列公式：
\[
\begin{aligned}
(x_1+x_2+\cdots+x_n)^2
={}&x_1^2+x_2^2+\cdots+x_n^2\\
&+2x_1x_2+2x_1x_3+\cdots+2x_1x_n\\
&+2x_2x_3+\cdots+2x_2x_n\\
&+\cdots\\
&+2x_{n-1}x_n.
\end{aligned}
\]
我们通过例子来介绍配方法.

\noindent\textbf{例 8.2.1}\quad 将下列二次型化成对角型：
\[
f(x_1,x_2,x_3)=x_1^2+2x_1x_2-4x_1x_3-3x_2^2-6x_2x_3+x_3^2.
\]

\noindent\textbf{解}\quad 先将含有 $x_1$ 的项放在一起凑成完全平方再减去必要的项：
\[
\begin{aligned}
f(x_1,x_2,x_3)
&=(x_1^2+2x_1x_2-4x_1x_3)-3x_2^2-6x_2x_3+x_3^2\\
&=\bigl((x_1+x_2-2x_3)^2-x_2^2-4x_3^2+4x_2x_3\bigr)-3x_2^2-6x_2x_3+x_3^2\\
&=(x_1+x_2-2x_3)^2-4x_2^2-2x_2x_3-3x_3^2.
\end{aligned}
\]
再对后面那些项配方：
\[
\begin{aligned}
-4x_2^2-2x_2x_3-3x_3^2
&=-\left(\left(2x_2+\frac12x_3\right)^2-\frac14x_3^2\right)-3x_3^2\\
&=-\left(2x_2+\frac12x_3\right)^2-\frac{11}{4}x_3^2.
\end{aligned}
\]
于是
\[
f(x_1,x_2,x_3)=(x_1+x_2-2x_3)^2-\left(2x_2+\frac12x_3\right)^2-\frac{11}{4}x_3^2.
\]
令
\[
\left\{\begin{array}{l}
y_1=x_1+x_2-2x_3,\\
y_2=2x_2+\dfrac12x_3,\\
y_3=x_3,
\end{array}\right.
\]
则
\[
\left(\begin{array}{c}
y_1\\
y_2\\
y_3
\end{array}\right)
=
\left(\begin{array}{ccc}
1&1&-2\\
0&2&\dfrac12\\
0&0&1
\end{array}\right)
\left(\begin{array}{c}
x_1\\
x_2\\
x_3
\end{array}\right),
\]
因此 $f=y_1^2-y_2^2-\dfrac{11}{4}y_3^2$，其变换矩阵为
\[
C=
\left(\begin{array}{ccc}
1&1&-2\\
0&2&\dfrac12\\
0&0&1
\end{array}\right)^{-1}
=
\left(\begin{array}{ccc}
1&-\dfrac12&\dfrac94\\
0&\dfrac12&-\dfrac14\\
0&0&1
\end{array}\right).
\]

\noindent\textbf{注}\quad 在用配方法化二次型为只含平方项的标准型的过程中，必须保证变换矩阵 $C$ 是非异阵. 如果我们按照例 8.2.1 的方法，将含 $x_1$ 的项放在一起配成一个完全平方，接下来将含 $x_2$ 的项放在一起再配方，如此不断做下去，最后得到的变换矩阵 $C$ 是一个主对角元全不为零的上三角阵，因此是一个非异阵. 有时我们用看似简单的方法得到的结果未必正确，比如用观察法即可得到下列配方：
\[
\begin{aligned}
f&=2x_1^2+2x_2^2+2x_3^2-2x_1x_2+2x_1x_3+2x_2x_3\\
&=(x_1-x_2)^2+(x_1+x_3)^2+(x_2+x_3)^2.
\end{aligned}
\]
若令 $y_1=x_1-x_2,\ y_2=x_1+x_3,\ y_3=x_2+x_3$，则 $f=y_1^2+y_2^2+y_3^2$. 由于矩阵
\[
\left(\begin{array}{ccc}
1&-1&0\\
1&0&1\\
0&1&1
\end{array}\right)
\]
不是非异阵，因此上述配方不是我们所需要的结论.

如果已知的二次型中没有平方项，我们可以采用下面例子中的方法.

\noindent\textbf{例 8.2.2}\quad 将二次型
\[
f(x_1,x_2,x_3,x_4)=2x_1x_2-x_1x_3+x_1x_4-x_2x_3+x_2x_4-2x_3x_4
\]
化成对角型.

\noindent\textbf{解}\quad 这个二次型缺少了 $x_i^2$ 项，因此无法用例 8.2.1 的方法配方，但我们可作如下变换：
\[
\left\{\begin{array}{l}
x_1=y_1+y_2,\\
x_2=y_1-y_2,\\
x_3=y_3,\\
x_4=y_4.
\end{array}\right.
\]
代入原二次型得
\[
f=2y_1^2-2y_2^2-2y_1y_3+2y_1y_4-2y_3y_4.
\]
这时 $y_1^2$ 项不为零，于是
\[
\begin{aligned}
f&=(2y_1^2-2y_1y_3+2y_1y_4)-2y_2^2-2y_3y_4\\
&=2\left(\left(y_1-\frac12y_3+\frac12y_4\right)^2-\frac14y_3^2-\frac14y_4^2+\frac12y_3y_4\right)-2y_2^2-2y_3y_4\\
&=2\left(y_1-\frac12y_3+\frac12y_4\right)^2-2y_2^2-\frac12y_3^2-y_3y_4-\frac12y_4^2\\
&=2\left(y_1-\frac12y_3+\frac12y_4\right)^2-2y_2^2-\frac12(y_3+y_4)^2.
\end{aligned}
\]
令
\[
\left\{\begin{array}{l}
z_1=y_1-\dfrac12y_3+\dfrac12y_4,\\
z_2=y_2,\\
z_3=y_3+y_4,\\
z_4=y_4,
\end{array}\right.
\]
于是
\[
f=2z_1^2-2z_2^2-\frac12z_3^2,
\]
其中 $z_4^2$ 的系数为零，故未写出.

为求变换矩阵 $C$，可从上面 $z_i$ 的表示式中解出 $y_i$：
\[
\left\{\begin{array}{l}
y_1=z_1+\dfrac12z_3-z_4,\\
y_2=z_2,\\
y_3=z_3-z_4,\\
y_4=z_4,
\end{array}\right.
\]
再将 $x_i$ 求出：
\[
\left\{\begin{array}{l}
x_1=z_1+z_2+\dfrac12z_3-z_4,\\
x_2=z_1-z_2+\dfrac12z_3-z_4,\\
x_3=z_3-z_4,\\
x_4=z_4,
\end{array}\right.
\]
于是
\[
C=\left(\begin{array}{cccc}
1&1&\dfrac12&-1\\
1&-1&\dfrac12&-1\\
0&0&1&-1\\
0&0&0&1
\end{array}\right).
\]

\subsection*{二、初等变换法}

用配方法化简二次型有时比较麻烦，求非异阵 $C$ 也比较麻烦. 我们常常用初等变换法来化简二次型，初等变换法的依据是引理 8.1.1.

下面我们通过例子来说明这种方法.

\noindent\textbf{例 8.2.3}\quad 将下列二次型化为对角型：
\[
f(x_1,x_2,x_3)=x_1^2-3x_2^2-2x_1x_2+2x_1x_3-6x_2x_3.
\]

\noindent\textbf{解}\quad 记与 $f$ 相伴的对称阵为 $A$，写出 $(A;I_3)$ 并作初等变换：
\[
(A;I_3)=
\left(\begin{array}{ccc|ccc}
1&-1&1&1&0&0\\
-1&-3&-3&0&1&0\\
1&-3&0&0&0&1
\end{array}\right)\rightarrow
\left(\begin{array}{ccc|ccc}
1&-1&1&1&0&0\\
0&-4&-2&1&1&0\\
1&-3&0&0&0&1
\end{array}\right)\rightarrow
\left(\begin{array}{ccc|ccc}
1&0&1&1&0&0\\
0&-4&-2&1&1&0\\
1&-2&0&0&0&1
\end{array}\right)
\]
\[
\rightarrow
\left(\begin{array}{ccc|ccc}
1&0&1&1&0&0\\
0&-4&-2&1&1&0\\
0&-2&-1&-1&0&1
\end{array}\right)\rightarrow
\left(\begin{array}{ccc|ccc}
1&0&0&1&0&0\\
0&-4&-2&1&1&0\\
0&-2&-1&-1&0&1
\end{array}\right)
\]
\[
\rightarrow
\left(\begin{array}{ccc|ccc}
1&0&0&1&0&0\\
0&-4&-2&1&1&0\\
0&0&0&-\dfrac32&-\dfrac12&1
\end{array}\right)\rightarrow
\left(\begin{array}{ccc|ccc}
1&0&0&1&0&0\\
0&-4&0&1&1&0\\
0&0&0&-\dfrac32&-\dfrac12&1
\end{array}\right).
\]
于是 $f$ 可化简为
\[
y_1^2-4y_2^2,
\]
\[
C=\left(\begin{array}{ccc}
1&0&0\\
1&1&0\\
-\dfrac32&-\dfrac12&1
\end{array}\right)'
=
\left(\begin{array}{ccc}
1&1&-\dfrac32\\
0&1&-\dfrac12\\
0&0&1
\end{array}\right).
\]

这种方法可总结如下：作 $n\times 2n$ 矩阵 $(A;I_n)$，对这个矩阵实施初等行变换，同时施以同样的初等列变换，将它左半边化为对角阵，则这个对角阵就是已化简的二次型的相伴矩阵，右半边的转置便是变换矩阵 $C$.

如碰到第 $(1,1)$ 元素是零的矩阵，可先设法将第 $(1,1)$ 元素化成非零，再进行上述过程.

\noindent\textbf{例 8.2.4}\quad 将二次型 $f(x_1,x_2,x_3)=2x_1x_2+4x_1x_3-4x_2x_3$ 化成对角型.

\noindent\textbf{解}\quad 写出与 $f$ 相伴的对称阵 $A$，作 $(A;I_3)$ 并将它的第二行加到第一行上，再将第二列加到第一列上：
\[
(A;I_3)=
\left(\begin{array}{ccc|ccc}
0&1&2&1&0&0\\
1&0&-2&0&1&0\\
2&-2&0&0&0&1
\end{array}\right)
\rightarrow
\left(\begin{array}{ccc|ccc}
2&1&0&1&1&0\\
1&0&-2&0&1&0\\
0&-2&0&0&0&1
\end{array}\right).
\]
同例 8.2.3 一样，对上述矩阵进行初等变换得到
\[
\left(\begin{array}{ccc|ccc}
2&0&0&1&1&0\\
0&-\dfrac12&0&-\dfrac12&\dfrac12&0\\
0&0&8&2&-2&1
\end{array}\right).
\]
因此 $f$ 化简为
\[
2y_1^2-\frac12y_2^2+8y_3^2,
\]
\[
C=\left(\begin{array}{ccc}
1&-\dfrac12&2\\
1&\dfrac12&-2\\
0&0&1
\end{array}\right).
\]

\section*{习题 8.2}

1. 用配方法把下列二次型化成对角型：

(1) $f(x_1,x_2,x_3)=x_1^2+x_2^2+3x_3^2+4x_1x_2+2x_1x_3+2x_2x_3$;

(2) $f(x_1,x_2,x_3)=x_1x_2+x_1x_3+x_2x_3$;

(3) $f(x_1,x_2,x_3,x_4)=x_1^2+x_2^2-2x_1x_2+4x_1x_3+2x_2x_3-2x_2x_4+2x_3x_4$;

(4) $f(x_1,x_2,x_3,x_4)=x_1x_2+x_2x_3+x_3x_4$.

2. 用初等变换法把下列二次型化为对角型并求出非异阵 $C$：

(1) $f(x_1,x_2,x_3)=x_1^2-3x_3^2-2x_1x_2+2x_1x_3-8x_2x_3$;

(2) $f(x_1,x_2,x_3)=2x_1x_2+2x_1x_3+2x_2x_3$;

(3) $f(x_1,x_2,x_3,x_4)=x_1^2+x_2^2+x_3^2+x_4^2+2x_1x_2+2x_1x_4-2x_2x_3-2x_2x_4-2x_3x_4$;

(4) $f(x_1,x_2,x_3,x_4)=2x_1x_2-2x_1x_4+2x_2x_3-4x_2x_4+6x_3x_4$.

3. 设矩阵
\[
A=\left(\begin{array}{cccc}
0&1&0&0\\
1&0&0&0\\
0&0&a&1\\
0&0&1&2
\end{array}\right)
\]
有一个特征值为 3，求 $a$ 的值并求可逆阵 $C$，使 $(AC)'(AC)$ 是对角阵.

\section*{§8.3\quad 惯性定理}

在上两节里我们已经知道如何将一个数域上的二次型化为标准型，即只含平方项的二次型. 在这一节里，我们将主要讨论实数域与复数域上的二次型.

\subsection*{一、实二次型}

我们已经知道，任意一个实对称阵 $A$ 必合同于一个对角阵：
\[
C'AC=\operatorname{diag}\{d_1,d_2,\cdots,d_r,0,\cdots,0\},
\]
其中 $d_i\neq0(i=1,\cdots,r)$. 注意到 $C$ 是可逆阵，故 $r=\operatorname{r}(C'AC)=\operatorname{r}(A)$，即秩 $r$ 是矩阵合同关系下的一个不变量. 如同相似标准型一样，我们的目的是要找出实对称阵在合同关系下的全系不变量，即找出一组足以判断两个实对称阵是否合同的合同不变量. 我们不妨设实对称阵已具有下列对角阵的形状：
\[
A=\operatorname{diag}\{d_1,d_2,\cdots,d_r,0,\cdots,0\}.
\]
由引理 8.1.1 不难知道，任意调换 $A$ 的主对角线上的元素得到的矩阵仍与 $A$ 合同. 因此我们可把零放在一起，把正项与负项放在一起，即可设 $d_1>0,\cdots,d_p>0;\ d_{p+1}<0,\cdots,d_r<0$. $A$ 所代表的二次型为
\begin{equation}
f(x_1,x_2,\cdots,x_n)=d_1x_1^2+d_2x_2^2+\cdots+d_rx_r^2.
\tag{8.3.1}
\end{equation}
令
\[
\left\{\begin{array}{l}
y_1=\sqrt{d_1}x_1,\ \cdots,\ y_p=\sqrt{d_p}x_p,\\
y_{p+1}=\sqrt{-d_{p+1}}x_{p+1},\ \cdots,\ y_r=\sqrt{-d_r}x_r,\\
y_j=x_j\quad (j=r+1,\cdots,n),
\end{array}\right.
\]
则 (8.3.1) 式变为
\begin{equation}
f=y_1^2+\cdots+y_p^2-y_{p+1}^2-\cdots-y_r^2.
\tag{8.3.2}
\end{equation}
这一事实等价于说 $A$ 合同于下列对角阵：
\begin{equation}
\operatorname{diag}\{1,\cdots,1;-1,\cdots,-1;0,\cdots,0\},
\tag{8.3.3}
\end{equation}
其中有 $p$ 个 $1$，$q$ 个 $-1$，$n-r$ 个零.

现在我们要证明 (8.3.2) 式中的数 $p$ 及 $q=r-p$ 是两个合同不变量.

\noindent\textbf{定理 8.3.1}\quad 设 $f(x_1,x_2,\cdots,x_n)$ 是一个 $n$ 元实二次型，且 $f$ 可化为两个标准型：
\[
c_1y_1^2+\cdots+c_py_p^2-c_{p+1}y_{p+1}^2-\cdots-c_ry_r^2,
\]
\[
d_1z_1^2+\cdots+d_kz_k^2-d_{k+1}z_{k+1}^2-\cdots-d_rz_r^2,
\]
其中 $c_i>0,d_i>0$，则必有 $p=k$.

\noindent\textbf{证明}\quad 用反证法，设 $p>k$. 由前面的说明不妨设 $c_i$ 及 $d_i$ 均为 1，因此
\begin{equation}
y_1^2+\cdots+y_p^2-y_{p+1}^2-\cdots-y_r^2
=
z_1^2+\cdots+z_k^2-z_{k+1}^2-\cdots-z_r^2.
\tag{8.3.4}
\end{equation}
又设
\[
x=By,\quad x=Cz,
\]
其中
\[
x=\left(\begin{array}{c}x_1\\x_2\\\vdots\\x_n\end{array}\right),\quad
y=\left(\begin{array}{c}y_1\\y_2\\\vdots\\y_n\end{array}\right),\quad
z=\left(\begin{array}{c}z_1\\z_2\\\vdots\\z_n\end{array}\right).
\]
于是 $z=C^{-1}By$. 令
\[
C^{-1}B=\left(\begin{array}{cccc}
c_{11}&c_{12}&\cdots&c_{1n}\\
c_{21}&c_{22}&\cdots&c_{2n}\\
\vdots&\vdots&&\vdots\\
c_{n1}&c_{n2}&\cdots&c_{nn}
\end{array}\right),
\]
则
\[
\left\{\begin{array}{l}
z_1=c_{11}y_1+c_{12}y_2+\cdots+c_{1n}y_n,\\
z_2=c_{21}y_1+c_{22}y_2+\cdots+c_{2n}y_n,\\
\cdots\cdots\cdots,\\
z_n=c_{n1}y_1+c_{n2}y_2+\cdots+c_{nn}y_n.
\end{array}\right.
\]
因为 $p>k$，故齐次线性方程组
\[
\left\{\begin{array}{l}
c_{11}y_1+c_{12}y_2+\cdots+c_{1n}y_n=0,\\
\cdots\cdots\cdots,\\
c_{k1}y_1+c_{k2}y_2+\cdots+c_{kn}y_n=0,\\
y_{p+1}=0,\\
\cdots\cdots\cdots,\\
y_n=0
\end{array}\right.
\]
必有非零解 ($n$ 个未知数，$n-(p-k)$ 个方程式). 令其中一个非零解为 $y_1=a_1,\cdots,y_p=a_p,y_{p+1}=0,\cdots,y_n=0$，把这组解代入 (8.3.4) 式左边得到
\[
a_1^2+\cdots+a_p^2>0.
\]
但这时 $z_1=\cdots=z_k=0$，故 (8.3.4) 式右边将小于等于零，引出了矛盾. 同理可证 $p<k$ 也不可能. $\square$

定理 8.3.1 通常称为惯性定理，(8.3.2) 式中的二次型称为 $f$ 的规范标准型.

\noindent\textbf{定义 8.3.1}\quad 设 $f(x_1,x_2,\cdots,x_n)$ 是一个实二次型，若它能化为形如 (8.3.2) 式的形状，则称 $r$ 是该二次型的秩，$p$ 是它的正惯性指数，$q=r-p$ 是它的负惯性指数，$s=p-q$ 称为 $f$ 的符号差.

显然，若已知秩 $r$ 与符号差 $s$，则 $p=\dfrac12(r+s)$，$q=\dfrac12(r-s)$. 事实上，在 $p,q,r,s$ 中只需知道其中两个数，其余两个数也就知道了. 由于实对称阵与实二次型之间的等价关系，我们将实二次型的秩、惯性指数及符号差也称为相应的实对称阵的秩、惯性指数及符号差.

\noindent\textbf{定理 8.3.2}\quad 秩与符号差 (或正负惯性指数) 是实对称阵在合同关系下的全系不变量.

\noindent\textbf{证明}\quad 由上面的定理知道，秩 $r$ 与符号差 $s$ 是实对称阵合同关系的不变量. 反之，若 $n$ 阶实对称阵 $A,B$ 的秩都为 $r$，符号差都是 $s$，则它们都合同于
\[
\operatorname{diag}\{1,\cdots,1;-1,\cdots,-1;0,\cdots,0\},
\]
其中有 $p=\dfrac12(r+s)$ 个 $1$，$q=\dfrac12(r-s)$ 个 $-1$ 及 $n-r$ 个零，因此 $A$ 与 $B$ 合同. 对正负惯性指数的结论也同样成立. $\square$

\subsection*{二、复二次型}

复二次型要比实二次型简单得多. 因为复二次型
\[
f(x_1,x_2,\cdots,x_n)=d_1x_1^2+d_2x_2^2+\cdots+d_rx_r^2
\]
必可化为
\[
z_1^2+z_2^2+\cdots+z_r^2,
\]
其中 $z_i=\sqrt{d_i}x_i(i=1,2,\cdots,r),\ z_j=x_j(j=r+1,\cdots,n)$. 所以复对称阵的合同关系只有一个全系不变量，那就是秩 $r$.

\section*{习题 8.3}

1. 把互相合同的实对称阵作为一个类，问：$n$ 阶实对称阵共有多少个类？

2. 举例说明实对称阵 $A$ 和它的伴随阵 $A^*$ 未必合同.

3. 设实二次型 $f$ 和 $g$ 的相伴矩阵分别是 $A$ 和 $A^{-1}$，求证：$f$ 和 $g$ 有相同的正负惯性指数.

4. 设实二次型 $f$ 的相伴矩阵为 $A$，若 $\det A<0$，求证：必存在一组实数 $a_1,a_2,\cdots,a_n$，使 $f(a_1,a_2,\cdots,a_n)<0$.

5. 设 $A$ 是 $m\times n$ 实矩阵且 $\operatorname{r}(A)=r$，作 $n$ 元二次型 $f(x)=x'(A'A)x$，求证：$f$ 的秩等于 $r$.

6. 设有 $n$ 元实二次型
\[
f(x_1,x_2,\cdots,x_n)=(x_1-a_1x_2)^2+(x_2-a_2x_3)^2+\cdots+(x_{n-1}-a_{n-1}x_n)^2+(x_n-a_nx_1)^2,
\]
其中 $a_i$ 都是实数，求证：$f$ 的正惯性指数大于等于 $n-1$，并确定 $a_i$ 适合什么条件时，$f$ 的正惯性指数等于 $n$.

7. 证明：一个秩大于 $1$ 的实二次型可以分解为两个实系数一次多项式之积的充分必要条件是它的秩等于 $2$ 且符号差等于零.

8. 设实二次型
\[
f(x_1,x_2,\cdots,x_n)=y_1^2+\cdots+y_k^2-y_{k+1}^2-\cdots-y_{k+s}^2,
\]
其中 $y_i=a_{i1}x_1+a_{i2}x_2+\cdots+a_{in}x_n\ (i=1,2,\cdots,k+s)$，求证：$f$ 的正惯性指数 $p\leq k$，负惯性指数 $q\leq s$.

9. 设分块实对称阵
\[
M=\left(\begin{array}{cc}A&O\\ O&B\end{array}\right),
\]
求证：$M$ 的正惯性指数等于 $A,B$ 的正惯性指数之和，$M$ 的负惯性指数等于 $A,B$ 的负惯性指数之和.

10. 设 $A$ 是 $n$ 阶实可逆阵，
\[
B=\left(\begin{array}{cc}O&A\\ A'&O\end{array}\right),
\]
求矩阵 $B$ 的正负惯性指数.

\section*{§8.4\quad 正定型与正定矩阵}

这一节里的二次型都假设是实二次型.

\noindent\textbf{定义 8.4.1}\quad 设 $f(x_1,x_2,\cdots,x_n)=x'Ax$ 是 $n$ 元实二次型，$A$ 是相伴矩阵.

(1) 若对任意 $n$ 维非零列向量 $\alpha$ 均有 $\alpha'A\alpha>0$，则称 $f$ 是正定二次型（简称正定型），矩阵 $A$ 称为正定矩阵（简称正定阵）；

(2) 若对任意 $n$ 维非零列向量 $\alpha$ 均有 $\alpha'A\alpha<0$，则称 $f$ 是负定二次型（简称负定型），矩阵 $A$ 称为负定矩阵（简称负定阵）；

(3) 若对任意 $n$ 维非零列向量 $\alpha$ 均有 $\alpha'A\alpha\geq 0$，则称 $f$ 是半正定二次型（简称半正定型），矩阵 $A$ 称为半正定矩阵（简称半正定阵）；

(4) 若对任意 $n$ 维非零列向量 $\alpha$ 均有 $\alpha'A\alpha\leq 0$，则称 $f$ 是半负定二次型（简称半负定型），矩阵 $A$ 称为半负定矩阵（简称半负定阵）；

(5) 若存在 $\alpha$，使 $\alpha'A\alpha>0$，又存在 $\beta$，使 $\beta'A\beta<0$，则称 $f$ 是不定型.

显然
\[
f(x_1,x_2,\cdots,x_n)=x_1^2+x_2^2+\cdots+x_n^2
\]
是正定型，而
\[
f(x_1,x_2,\cdots,x_n)=-x_1^2-x_2^2-\cdots-x_n^2
\]
是负定型. 事实上，我们可以证明如下的定理.

\noindent\textbf{定理 8.4.1}\quad 设 $f(x_1,x_2,\cdots,x_n)$ 是 $n$ 元实二次型，则 $f$ 是正定型的充分必要条件是 $f$ 的正惯性指数等于 $n$；$f$ 是负定型的充分必要条件是 $f$ 的负惯性指数等于 $n$；$f$ 是半正定型的充分必要条件是 $f$ 的正惯性指数等于 $f$ 的秩 $r$；$f$ 是半负定型的充分必要条件是 $f$ 的负惯性指数等于 $f$ 的秩 $r$.

\noindent\textbf{证明}\quad 若 $f$ 的正惯性指数等于 $n$，则 $f$ 可化为下列标准型：
\[
f=y_1^2+y_2^2+\cdots+y_n^2,
\]
显然 $f$ 是正定型. 反之，若 $f$ 是正定型，如果 $f$ 的正惯性指数 $p<n$，则 $f$ 可化为如下标准型：
\begin{equation}
f=y_1^2+\cdots+y_p^2-c_{p+1}y_{p+1}^2-\cdots-c_ny_n^2,
\tag{8.4.1}
\end{equation}
其中 $c_j\geq 0(j=p+1,\cdots,n)$. 这时令 $b_1=\cdots=b_p=0,\ b_{p+1}=\cdots=b_n=1$，则 $b_1,b_2,\cdots,b_n$ 不全为零. 假设这时 $x=Cy$，其中
\[
x=\left(\begin{array}{c}x_1\\x_2\\ \vdots\\x_n\end{array}\right),\qquad
y=\left(\begin{array}{c}y_1\\y_2\\ \vdots\\y_n\end{array}\right),
\]
$C$ 是非异阵，则从 $y_i=b_i(i=1,\cdots,n)$ 可得 $x_i=a_i(i=1,\cdots,n)$ 是一组不全为零的实数，于是
\[
f(a_1,a_2,\cdots,a_n)\leq 0,
\]
这与 $f$ 是正定型矛盾. 其余结论的证明类似. $\square$

显然我们有下面的定理.

\noindent\textbf{定理 8.4.2}\quad $n$ 阶实对称阵 $A$ 是正定阵当且仅当它合同于单位阵 $I_n$；$A$ 是负定阵当且仅当它合同于 $-I_n$；$A$ 是半正定阵当且仅当 $A$ 合同于下列对角阵：
\[
\left(\begin{array}{cc}I_r&O\\ O&O\end{array}\right);
\]
$A$ 是半负定阵当且仅当 $A$ 合同于下列对角阵：
\[
\left(\begin{array}{cc}-I_r&O\\ O&O\end{array}\right).
\]

\noindent\textbf{定义 8.4.2}\quad 设 $A=(a_{ij})$ 是 $n$ 阶矩阵，$A$ 的 $n$ 个子式：
\[
\left|\begin{array}{cccc}
a_{11}&a_{12}&\cdots&a_{1k}\\
a_{21}&a_{22}&\cdots&a_{2k}\\
\vdots&\vdots&&\vdots\\
a_{k1}&a_{k2}&\cdots&a_{kk}
\end{array}\right|\quad (k=1,2,\cdots,n)
\]
称为 $A$ 的顺序主子式.

\noindent\textbf{定理 8.4.3}\quad $n$ 阶实对称阵 $A$ 是正定阵的充分必要条件是它的 $n$ 个顺序主子式全大于零.

\noindent\textbf{证明}\quad 先证必要性. 设 $n$ 阶实对称阵 $A=(a_{ij})$ 为正定阵，则对应的实二次型
\[
f(x_1,x_2,\cdots,x_n)=\sum_{j=1}^n\sum_{i=1}^n a_{ij}x_ix_j
\]
为正定型. 令
\[
f_k(x_1,x_2,\cdots,x_k)=\sum_{j=1}^k\sum_{i=1}^k a_{ij}x_ix_j,
\]
则对任意一组不全为零的实数 $c_1,c_2,\cdots,c_k$，有
\[
f_k(c_1,c_2,\cdots,c_k)=f(c_1,c_2,\cdots,c_k,0,\cdots,0)>0,
\]
因此 $f_k$ 是一个正定二次型，从而它的相伴矩阵 $A_k$（由 $A$ 的前 $k$ 行及前 $k$ 列组成）是一个正定阵. 由于 $A_k$ 合同于 $I_k$，故存在 $k$ 阶非异阵 $B$，使
\[
B'A_kB=I_k,
\]
于是
\[
\det(B'A_kB)=\det(B)^2\det(A_k)=1,
\]
即有 $\det(A_k)>0(k=1,2,\cdots,n)$.

再证充分性. 对 $A$ 的阶数进行归纳. 当 $n=1$ 时，$A=(a),\ a>0$，于是 $f=ax_1^2$ 是正定型，从而 $A$ 是正定阵. 设结论对 $n-1$ 成立，现证明对 $n$ 阶实对称阵 $A$，若它的 $n$ 个顺序主子式全大于零，则 $A$ 必是正定阵. 记 $A_{n-1}$ 是 $A$ 的 $n-1$ 阶顺序主子式所在的矩阵，则 $A$ 可写为
\[
A=\left(\begin{array}{cc}A_{n-1}&\alpha\\ \alpha'&a_{nn}\end{array}\right).
\]
因为 $A$ 的顺序主子式全大于零，故 $A_{n-1}$ 的顺序主子式也全大于零，由归纳假设，$A_{n-1}$ 是正定阵. 于是 $A_{n-1}$ 合同于 $n-1$ 阶单位阵，即存在 $n-1$ 阶非异阵 $B$，使
\[
B'A_{n-1}B=I_{n-1}.
\]
令 $C$ 是下列分块矩阵：
\[
C=\left(\begin{array}{cc}B&O\\ O&1\end{array}\right),
\]
则
\[
C'AC=\left(\begin{array}{cc}B'&O\\ O&1\end{array}\right)
\left(\begin{array}{cc}A_{n-1}&\alpha\\ \alpha'&a_{nn}\end{array}\right)
\left(\begin{array}{cc}B&O\\ O&1\end{array}\right)
=\left(\begin{array}{cc}I_{n-1}&B'\alpha\\ \alpha'B&a_{nn}\end{array}\right).
\]
这是一个实对称阵，其形式为
\[
C'AC=\left(\begin{array}{ccccc}
1&\cdots&0&c_1\\
\vdots&&\vdots&\vdots\\
0&\cdots&1&c_{n-1}\\
c_1&\cdots&c_{n-1}&a_{nn}
\end{array}\right),
\]
用第三类初等行及列变换可将上述矩阵化为对角阵. 这相当于对 $C'AC$ 右乘一个非异阵 $Q$ 后，再左乘 $Q'$ 得到一个对角阵，亦即 $Q'C'ACQ$ 等于
\[
\operatorname{diag}\{1,\cdots,1,c\}.
\]
由于 $|A|>0$，故 $c>0$，这就证明了 $A$ 是一个正定阵. $\square$

\noindent\textbf{例 8.4.1}\quad 试求 $t$ 的取值范围，使下列二次型为正定型：
\[
f(x_1,x_2,x_3,x_4)=x_1^2+4x_2^2+4x_3^2+3x_4^2+2tx_1x_2-2x_1x_3+4x_2x_3.
\]

\noindent\textbf{解}\quad 这个二次型的相伴矩阵为
\[
A=\left(\begin{array}{rrrr}
1&t&-1&0\\
t&4&2&0\\
-1&2&4&0\\
0&0&0&3
\end{array}\right).
\]
$A$ 的顺序主子式为
\[
|A_1|=1>0,\qquad |A_2|=\left|\begin{array}{cc}1&t\\ t&4\end{array}\right|=4-t^2,
\]
\[
|A_3|=\left|\begin{array}{rrr}
1&t&-1\\
t&4&2\\
-1&2&4
\end{array}\right|=-4(t-1)(t+2),\qquad |A_4|=-12(t-1)(t+2).
\]
要使 $A$ 正定，必须
\[
4-t^2>0,\qquad -4(t-1)(t+2)>0,
\]
得 $-2<t<1$.

\noindent\textbf{例 8.4.2}\quad 若 $A$ 是正定阵，证明：

(1) $A$ 的任一 $k$ 阶主子阵，即由 $A$ 的第 $i_1,i_2,\cdots,i_k$ 行及 $A$ 的第 $i_1,i_2,\cdots,i_k$ 列交点上元素组成的矩阵，必是正定阵；

(2) $A$ 的所有主子式全大于零，特别，$A$ 的主对角元素全大于零；

(3) $A$ 中绝对值最大的元素仅在主对角线上.

\noindent\textbf{证明}\quad (1) 设 $A_k$ 是矩阵 $A$ 的第 $k$ 个顺序主子式所在的矩阵，则 $A_k$ 是实对称阵且其顺序主子式都大于零，因此 $A_k$ 是正定阵.

经过若干次行对换以及相同的列对换，我们不难将 $A$ 的第 $i_1,i_2,\cdots,i_k$ 行及 $i_1,i_2,\cdots,i_k$ 列分别换成第 $1,2,\cdots,k$ 行和第 $1,2,\cdots,k$ 列，利用上面的结论即知 (1) 成立.

(2) 是 (1) 的推论.

(3) 用反证法. 假设 $a_{ij}(i\ne j)$ 是 $A$ 的绝对值最大的元素. 根据 (1)，我们只需证明由第 $i,j$ 行与第 $i,j$ 列交点上元素组成的矩阵不是正定阵即可. 考虑矩阵（由 $A$ 的对称性不妨设 $i<j$）
\[
\left(\begin{array}{cc}a_{ii}&a_{ij}\\ a_{ji}&a_{jj}\end{array}\right),
\]
注意到 $a_{ij}=a_{ji}$ 且 $|a_{ij}|\geq a_{ii}, |a_{ij}|\geq a_{jj}$，上述矩阵的行列式值 $a_{ii}a_{jj}-a_{ij}^2\leq 0$，所以这个矩阵一定不是正定阵. $\square$

\section*{习题 8.4}

1. 判定下列实二次型属于哪种型：

(1) $f(x_1,x_2,x_3)=x_1^2+2x_2^2+x_3^2+4x_1x_2-8x_1x_3+4x_2x_3$；

(2) $f(x_1,x_2,x_3)=x_1^2+4x_2^2+5x_3^2+4x_1x_2-4x_1x_3-8x_2x_3$.

2. 设下列实二次型是正定型，求 $\lambda$ 的取值范围：

(1) $f(x_1,x_2,x_3)=5x_1^2+x_2^2+\lambda x_3^2+4x_1x_2-2x_1x_3-2x_2x_3$；

(2) $f(x_1,x_2,x_3)=2x_1^2+x_2^2+3x_3^2+2\lambda x_1x_2+2x_1x_3$.

3. 设 $A$ 是 $n$ 阶实对称阵，$P_1,P_2,\cdots,P_n$ 是 $A$ 的 $n$ 个顺序主子式，求证：$A$ 是负定阵的充分必要条件是
\[
P_1<0,\quad P_2>0,\quad \cdots,\quad (-1)^nP_n>0.
\]

4. 求证：$n$ 阶实对称阵 $A$ 是正定阵的充分必要条件是 $A$ 的前 $n-1$ 个顺序主子式的代数余子式以及第 $n$ 个顺序主子式全大于零.

5. 设 $A$ 是 $m$ 阶正定实对称阵，$B$ 是 $m\times n$ 实矩阵. 求证：$B'AB$ 是正定阵的充分必要条件是 $\operatorname{r}(B)=n$.

6. 设实二次型 $f(x_1,x_2,\cdots,x_n)$ 仅在 $x_1=x_2=\cdots=x_n=0$ 时为零，求证：$f$ 必是正定型或负定型.

7. 若 $A,B$ 都是 $n$ 阶正定实对称阵，求证：$A^{-1}$ 与 $A+B$ 也都是正定阵.

8. 若 $A$ 是可逆半正定实对称阵，求证：$A$ 必是正定阵.

9. 设 $A$ 是半正定实对称阵，求证：$A$ 的任一主子阵都是半正定阵.

10. 设 $A$ 是 $n$ 阶实对称阵，求证：$A$ 是秩为 $r$ 的半正定阵的充分必要条件是存在秩为 $r$ 的 $r\times n$ 实矩阵 $B$，使 $A=B'B$.

11. 设 $A$ 是 $n$ 阶实对称阵，求证：$A$ 是半正定阵的充分必要条件是对任意的正实数 $c$，$A+cI_n$ 都是正定阵.

12. 设 $A$ 是 $n$ 阶实对称阵，证明：$A$ 是半正定阵的充分必要条件是 $A$ 的所有主子式都大于等于零. 举例说明 $A$ 的所有顺序主子式都大于等于零并不能保证 $A$ 是半正定阵.

13. 设 $A$ 为 $n$ 阶实对称阵，证明下列结论成立：

(1) 若 $A$ 是正定阵，则 $A^*$ 也是正定阵.

(2) 若 $A$ 是半正定阵，则 $A^*$ 也是半正定阵.

(3) 若 $A$ 是负定阵，则当 $n$ 是偶数时，$A^*$ 为负定阵；当 $n$ 是奇数时，$A^*$ 为正定阵.

14. 设 $A$ 是 $n$ 阶正定实对称阵，求证：
\[
B=\left(\begin{array}{cc}A&-I_n\\ -I_n&A^{-1}\end{array}\right)
\]
是半正定阵.

\section*{*§8.5\quad Hermite 型}

现在我们来考虑复数域上的 Hermite 型. 为方便起见，我们把 $n$ 个变元的 Hermite 型看成是复数域上的二次齐次函数：
\begin{equation}
f(x_1,x_2,\cdots,x_n)=\sum_{j=1}^n\sum_{i=1}^n a_{ij}\overline{x_i}x_j,
\tag{8.5.1}
\end{equation}
其中 $\overline{a_{ij}}=a_{ji}$. Hermite 型虽然系数是复数且变元 $x_i$ 是复数域上的变元，但作为函数它的值却总是实数，这点从 Hermite 型的定义即可看出. 事实上，
\[
\overline{f}=\sum_{j=1}^n\sum_{i=1}^n \overline{a_{ij}}x_i\overline{x_j}
=\sum_{i=1}^n\sum_{j=1}^n a_{ji}\overline{x_j}x_i=f,
\]
因此 $f$ 的值总是实数. 当 Hermite 型的变元 $x_i$ 取实变元且 $a_{ij}$ 都是实数时，$f$ 就是实二次型. 因此，Hermite 型也可看成是实二次型的推广. 事实上，Hermite 型有许多与实二次型相同的性质.

Hermite 型 (8.5.1) 可写成如下的矩阵相乘的形式：
\[
f(x_1,x_2,\cdots,x_n)=\overline{x}'Ax,
\]
其中
\[
A=\left(\begin{array}{cccc}
a_{11}&a_{12}&\cdots&a_{1n}\\
a_{21}&a_{22}&\cdots&a_{2n}\\
\vdots&\vdots&&\vdots\\
a_{n1}&a_{n2}&\cdots&a_{nn}
\end{array}\right),\qquad
x=\left(\begin{array}{c}x_1\\x_2\\ \vdots\\x_n\end{array}\right),
\]
且满足 $\overline{A}'=A$，这样的矩阵称为 Hermite 矩阵. 与实二次型类似，Hermite 型与 Hermite 矩阵之间有着一一对应关系，即给定一个 $n$ 变元的 Hermite 型必相伴有一个 $n$ 阶 Hermite 矩阵，反之给定一个 $n$ 阶 Hermite 矩阵，必有一个 $n$ 元 Hermite 型与之对应.

设 $x=Cy$，其中 $C$ 是一个非异复矩阵，$y=(y_1,y_2,\cdots,y_n)'$，则
\[
f=(\overline{Cy})'A(Cy)=\overline{y}'\,\overline{C}'ACy.
\]
矩阵 $\overline{C}'AC$ 仍是一个 Hermite 矩阵.

\noindent\textbf{定义 8.5.1}\quad 设 $A,B$ 是两个 Hermite 矩阵，若存在非异复矩阵 $C$，使
\[
B=\overline{C}'AC,
\]
则称 $A$ 与 $B$ 是复相合的.

容易证明复相合是一个等价关系. 与实二次型类似，Hermite 型
\begin{equation}
a_1\overline{y_1}y_1+a_2\overline{y_2}y_2+\cdots+a_n\overline{y_n}y_n
\tag{8.5.2}
\end{equation}
称为标准型. Hermite 型的基本问题是如何把一个 Hermite 型化成像 (8.5.2) 式那样的标准型. 这个问题等价于寻找非异阵 $C$，使 $\overline{C}'AC$ 成为对角阵. 类似于对称性，我们可以证明如下定理.

\noindent\textbf{定理 8.5.1}\quad 设 $A$ 是一个 Hermite 矩阵，则必存在一个非异阵 $C$，使 $\overline{C}'AC$ 是一个对角阵且主对角线上的元素都是实数.

\noindent\textbf{证明}\quad 寻找 $C$ 的过程类似于对称阵的情形，故从略. 现只需说明后面的结论. 事实上，$\overline{C}'AC$ 仍是 Hermite 矩阵，因此主对角线上的元素 $b_{ii}=\overline{b_{ii}}$，必是实数. $\square$

类似实二次型，我们可以证明 Hermite 型的惯性定理.

\noindent\textbf{定理 8.5.2}\quad 设 $f(x_1,x_2,\cdots,x_n)$ 是一个 Hermite 型，则它总可以化为如下标准型：
\begin{equation}
\overline{y_1}y_1+\cdots+\overline{y_p}y_p-\overline{y_{p+1}}y_{p+1}-\cdots-\overline{y_r}y_r,
\tag{8.5.3}
\end{equation}
且若 $f$ 又可化为另一个标准型：
\[
\overline{z_1}z_1+\cdots+\overline{z_k}z_k-\overline{z_{k+1}}z_{k+1}-\cdots-\overline{z_r}z_r,
\]
则 $p=k$.

我们称 (8.5.3) 式中的 $r$ 为 $f$ 的秩，$p$ 为 $f$ 的正惯性指数，$q=r-p$ 为 $f$ 的负惯性指数，$p-q$ 为 $f$ 的符号差. 同样地，秩与符号差（或正负惯性指数）是 Hermite 矩阵复相合的全系不变量. 上述这些结论的证明与实二次型是平行的，我们略去其证明，把它们留给读者作为练习.

由于 Hermite 型的值总是实数，故可定义正定型、负定型、半正定型、半负定型及不定型的概念. 相应地，有正定 Hermite 矩阵、负定 Hermite 矩阵、半正定 Hermite 矩阵和半负定 Hermite 矩阵的概念. 我们只对正定 Hermite 型和正定 Hermite 矩阵叙述其概念以及判定正定 Hermite 矩阵的定理.

\noindent\textbf{定义 8.5.2}\quad 设 $f(x_1,x_2,\cdots,x_n)$ 是 Hermite 型，若对任一组不全为零的复数 $c_1,c_2,\cdots,c_n$，均有
\[
f(c_1,c_2,\cdots,c_n)>0,
\]
则称 $f$ 是正定 Hermite 型，它对应的矩阵称为正定 Hermite 矩阵.

\noindent\textbf{定理 8.5.3}\quad $n$ 阶 Hermite 矩阵 $A$ 为正定的充分必要条件是它的 $n$ 个顺序主子式全大于零.

这个定理的证明类似于实对称阵的情形，故从略.

\section*{习题 8.5}

1. 求证：Hermite 矩阵 $A$ 的下列变换都是复相合变换：

(1) 对换 $A$ 的第 $i$ 行与第 $j$ 行，再对换第 $i$ 列与第 $j$ 列；

(2) 将非零复数 $k$ 乘以 $A$ 的第 $i$ 行，再将其轭复数 $\overline{k}$ 乘以第 $i$ 列；

(3) 将复数 $k$ 乘以 $A$ 的第 $i$ 行加到第 $j$ 行上，再将其轭复数 $\overline{k}$ 乘以第 $i$ 列加到第 $j$ 列上.

2. 利用上述复相合变换证明定理 8.5.1.

3. 仿照定理 8.3.1 的证明方法证明定理 8.5.2.

4. 设 $A$ 是 $n$ 阶 Hermite 矩阵，仿照定理 8.4.1 的证明方法证明下列命题等价：

(1) $A$ 是正定 Hermite 矩阵；

(2) $A$ 复相合于单位阵 $I_n$；

(3) 存在 $n$ 阶非异复方阵 $B$，使 $A=\overline{B}'B$.

5. 仿照定理 8.4.3 的证明方法证明定理 8.5.3.

\section*{历史与展望}

对特定二次型的研究，例如一个整数能否是一个整系数二次型的取值等问题，可以追溯到许多世纪以前. 一个著名的例子是关于两个整数平方和的 Fermat（费马）定理，它决定了一个整数何时可用 $x^2+y^2$ 的形式表示，其中 $x,y$ 都是整数. 1801 年，Gauss 出版了著作《算术研究》，建立了整系数二元二次型的算术理论，另外还引入了二次型的正定性和半正定性等重要的概念. 从那时起，二次型在数论之外的其他数学领域被进一步推广和研究.

将二次曲线和二次曲面的方程变形，选取主轴方向作为坐标轴以简化方程的形状，这个问题是在 18 世纪提出的. 当二次曲面的方程是标准型（即主轴是坐标轴）时，二次曲面用标准型平方项系数的符号来进行分类，这是 Cauchy 在 1826 年的论文《几何中无穷小演算的应用教程》中给出的. 然而那时并不清楚，在化简成标准型时，总得到相同数目的正项和负项. Sylvester 在 1852 年回答了这个问题，这就是 $n$ 元实二次型的惯性定理，但他认为这个定理是自明的，没有给出证明，直到 1857 年 Jacobi 重新发现和证明了这个定理.

二次型化简的进一步研究涉及二次型的特征方程和特征根的概念（其实就是二次型相伴实对称阵的特征多项式和特征值），一旦计算出二次型的特征根，就能立刻得到主轴的长度（参考定理 9.5.4）. 特征方程的概念隐含地出现在 Euler 的化三元二次型到它们的主轴上去的著作中，虽然他对特征根的实性缺乏证明. 特征方程的概念首先明确地出现在 Lagrange 关于线性微分方程组的著作中，也出现在 Laplace 在同一领域的著作中.

Cauchy 和 Weierstrass 还研究了二次型束 $uA+vB$ 的化简问题，其中 $A,B$ 是任意给定的两个二次型. 他们证明了：若 $A,B$ 中有一个是正定型，则它们可同时化简为标准型. Sylvester 在研究二次曲面束 $A+\lambda B$ 的分类问题中，引入了行列式因子、不变因子和初等因子等概念. 这些关于二次型的研究极大促进了行列式理论和矩阵理论的发展.

\section*{复习题八}

1. 设 $A$ 为 $n$ 阶实对称阵，证明下列命题等价：

(1) $A$ 是正定阵；

(2) 存在主对角元全为 $1$ 的上三角阵 $B$，使 $A=B'DB$，其中 $D$ 是主对角元全为正数的对角阵；

(3) 存在主对角元全为正数的上三角阵 $C$，使 $A=C'C$.

注：正定实对称阵 $A$ 的上述分解 (2) 和 (3) 都是存在并且唯一的，即满足上述等式的矩阵 $B,C,D$ 是存在并且唯一的. 分解 (3) 通常称为正定实对称阵 $A$ 的 Cholesky（楚列斯基）分解.

2. 设 $A$ 是数域 $\mathbb K$ 上的 $n$ 阶非异阵，若存在主对角元全为 $1$ 的下三角阵 $L$ 以及上三角阵 $U$，使 $A=LU$，则称这一分解为矩阵 $A$ 的 $LU$ 分解（$L$ 表示下三角，$U$ 表示上三角）. 证明：$n$ 阶非异阵 $A$ 存在 $LU$ 分解的充分必要条件是 $A$ 的 $n$ 个顺序主子式都不等于零. 此时，$A$ 的 $LU$ 分解是唯一的.

3. 设 $f(x)=x'Ax$ 是实二次型，相伴矩阵 $A$ 的前 $n-1$ 个顺序主子式 $P_1,\cdots,P_{n-1}$ 非零，求证：经过可逆线性变换 $f$ 可化为下列标准型：
\[
f=P_1y_1^2+\frac{P_2}{P_1}y_2^2+\cdots+\frac{P_n}{P_{n-1}}y_n^2,
\]
其中 $P_n=|A|$.

4. 设 $A$ 为 $n$ 阶正定实对称阵且非主对角元都是负数，求证：$A^{-1}$ 的每个元素都是正数.

5. 设 $A=(a_{ij})$ 是 $n$ 阶正定实对称阵，其逆阵 $A^{-1}=(b_{ij})$，求证：$a_{ii}b_{ii}\geq 1$，且等号成立当且仅当 $A$ 的第 $i$ 行和第 $i$ 列的所有元素除了 $a_{ii}$ 之外全为零.

6. 设分块实对称阵
\[
M=\left(\begin{array}{cc}A&C\\ C'&B\end{array}\right),
\]
其中 $A,B$ 都可逆，用 $p(A),q(A)$ 分别表示 $A$ 的正
负惯性指数，求证：
\[
p(A)+p(B-C'A^{-1}C)=p(B)+p(A-CB^{-1}C'),
\]
\[
q(A)+q(B-C'A^{-1}C)=q(B)+q(A-CB^{-1}C').
\]

7. 设 $\alpha$ 是 $n$ 维实列向量且 $\alpha'\alpha=1$，求矩阵 $I_n-2\alpha\alpha'$ 的正负惯性指数.

8. 求 $n(n\geq 2)$ 阶实对称阵 $A$ 的正负惯性指数，其中 $a_i$ 均为实数：
\[
A=\left(\begin{array}{cccc}
a_1^2&a_1a_2+1&\cdots&a_1a_n+1\\
a_2a_1+1&a_2^2&\cdots&a_2a_n+1\\
\vdots&\vdots&&\vdots\\
a_na_1+1&a_na_2+1&\cdots&a_n^2
\end{array}\right).
\]

9. 求证：任一 $n$ 阶复矩阵 $A$ 都相似于一个复对称阵.

10. 设 $A$ 为 $n$ 阶正定实对称阵，$\alpha,\beta$ 为 $n$ 维实列向量，证明：
\[
\alpha'A\alpha+\beta'A^{-1}\beta\geq 2\alpha'\beta,
\]
且等号成立的充分必要条件是 $A\alpha=\beta$.

11. 设 $A$ 为 $n$ 阶正定实对称阵，$\alpha,\beta$ 为 $n$ 维实列向量，证明：
\[
(\alpha'\beta)^2\leq(\alpha'A\alpha)(\beta'A^{-1}\beta),
\]
且等号成立的充分必要条件是 $A\alpha$ 与 $\beta$ 成比例.

12. 设 $A$ 为 $n$ 阶实对称阵，证明：

(1) 若 $A$ 可逆，则 $A$ 为正定阵的充分必要条件是对任意的 $n$ 阶正定实对称阵 $B$，$\operatorname{tr}(AB)>0$；

(2) $A$ 为半正定阵的充分必要条件是对任意的 $n$ 阶半正定实对称阵 $B$，$\operatorname{tr}(AB)\geq 0$.

13. 设 $A,B$ 都是 $n$ 阶半正定实对称阵，证明：$AB=O$ 的充分必要条件是 $\operatorname{tr}(AB)=0$.

14. 化下列实二次型为标准型：
\[
f(x_1,x_2,\cdots,x_n)=x_1x_2+x_2x_3+\cdots+x_{n-1}x_n.
\]

15. 化下列实二次型为标准型：
\[
f(x_1,x_2,\cdots,x_n)=\sum_{i=1}^n x_i^2+\sum_{1\leq i<j\leq n}x_ix_j.
\]

16. 化下列实二次型为标准型：
\[
f(x_1,x_2,\cdots,x_n)=\sum_{i=1}^n(x_i-s)^2,\qquad s=\frac1n(x_1+x_2+\cdots+x_n).
\]

17. 设 $X=(x_{ij})_{n\times n}$ 是 $n$ 阶矩阵变量，$f(X)=\operatorname{tr}(X^2)$ 是关于未定元 $x_{ij}(i,j=1,\cdots,n)$ 的实二次型，试求 $f$ 的正负惯性指数.

18. 设 $A$ 为 $m$ 阶实对称阵，$C$ 为 $m\times n$ 实矩阵，证明：$C'AC$ 的正惯性指数小于等于 $A$ 的正惯性指数；$C'AC$ 的负惯性指数小于等于 $A$ 的负惯性指数.

19. 设 $A,B$ 为 $n$ 阶实对称阵，并用 $p(A),q(A)$ 分别表示 $A$ 的正负惯性指数. 求证：
\[
p(A+B)\leq p(A)+p(B),\qquad q(A+B)\leq q(A)+q(B).
\]

20. 设 $A$ 是 $n$ 阶正定实对称阵，求证：函数
\[
f(x)=x'Ax+2\beta'x+c
\]
的极小值等于 $c-\beta'A^{-1}\beta$，其中 $\beta=(b_1,\cdots,b_n)'$，$b_i$ 和 $c$ 都是实数.

21. 求下列实二次型的标准型：

(1) $f(x_1,x_2,\cdots,x_n)=\displaystyle\sum_{i,j=1}^n \max\{i,j\}x_ix_j$；

(2) $f(x_1,x_2,\cdots,x_n)=\displaystyle\sum_{i,j=1}^n |i-j|x_ix_j$.

22. 设 $A=(a_{ij}),B=(b_{ij})$ 都是 $n$ 阶正定实对称阵，求证：$H=(a_{ij}b_{ij})$ 也是正定阵.

23. 设 $A$ 是 $n$ 阶可逆实对称阵，$S$ 是 $n$ 阶实反对称阵且 $AS=SA$，求证：$A+S$ 是可逆阵.

24. 设 $A$ 是 $n$ 阶正定实对称阵，$S$ 是 $n$ 阶实反对称阵，求证：

(1) $|A+S|\geq |A|+|S|$，且等号成立当且仅当 $n\leq 2$ 或当 $n\geq 3$ 时，$S=O$.

(2) $|A+S|\geq |A|$，且等号成立当且仅当 $S=O$.

25. 设 $A,B,A-B$ 都是 $n$ 阶正定实对称阵，求证：$B^{-1}-A^{-1}$ 也是正定阵.

26. 设 $A$ 为 $n$ 阶正定实对称阵，$n$ 维实列向量 $\alpha,\beta$ 满足 $\alpha'\beta>0$，求证：
\[
H=A-\frac{A\beta\beta'A}{\beta'A\beta}+\frac{\alpha\alpha'}{\alpha'\beta}
\]
是正定阵.

27. 求证：$n$ 阶实对称阵 $A=(a_{ij})$ 是正定阵，其中
\[
a_{ij}=\frac{1}{i+j}.
\]

28. 设 $A$ 是 $n$ 阶实对称阵，求证：若 $A$ 是主对角元全大于零的严格对角占优阵，则 $A$ 是正定阵.

29. 设 $A$ 是 $n$ 阶实对称阵，求证：必存在正实数 $k$，使对任一 $n$ 维实列向量 $\alpha$，总有
\[
-k\alpha'\alpha\leq \alpha'A\alpha\leq k\alpha'\alpha.
\]

30. 设 $\alpha,\beta$ 为 $n$ 维非零实列向量，求证：$\alpha'\beta>0$ 成立的充要条件是存在 $n$ 阶正定实对称阵 $A$，使得 $\alpha=A\beta$.

31. 设 $A,B$ 是 $n$ 阶实矩阵，使得 $A'B'+BA$ 是正定阵，求证：$A,B$ 都是非异阵.

32. 设 $A,B,C$ 都是 $n$ 阶正定实对称阵，$g(t)=|t^2A+tB+C|$ 是关于 $t$ 的多项式，求证：$g(t)$ 所有复根的实部都小于零.

33. 设 $A=(a_{ij})$ 是 $n$ 阶正定实对称阵，$P_{n-1}$ 是 $A$ 的第 $n-1$ 个顺序主子式，求证：
\[
|A|\leq a_{nn}P_{n-1}.
\]

34. 设 $A=(a_{ij})$ 是 $n$ 阶正定实对称阵，求证：
\[
|A|\leq a_{11}a_{22}\cdots a_{nn},
\]
且等号成立当且仅当 $A$ 是对角阵.

\end{document}
